% !PRINTED 40
\weberpage{40}
Denn sind $\alpha$ und $\beta$ beliebig gegeben, so kann man $s$ nach dem Modul $ab$ so bestimmen, dass
\[
s \equiv \alpha \ (\mathrm{mod}\ a), \quad s \equiv \beta \ (\mathrm{mod}\ b),
\]
wodurch (2) mit (1) identisch wird. Trotzdem ist in gewissem Sinne die Constitution der Basis durch die Natur der Gruppen völlig bestimmt, nach dem folgenden Satze:

\gesperrt{II. Sind
\[
A_1, A_2 \ldots A_\nu
\]
mit den Graden
\[
a_1, a_2 \ldots a_\nu
\]
und
\[
B_1, B_2 \ldots B_\mu
\]
mit den Graden
\[
b_1, b_2 \ldots b_\mu
\]
zwei Basen einer Abel'schen Gruppe $S$ vom Grade $n$, ist $p$ eine in $n$ aufgehende Primzahl und
\weq{1}{a_1 = p_1 a_1', \quad a_2 = p_2 a_2', \quad \ldots \quad a_\nu = p_\nu a_\nu',}
\weq{2}{b_1 = p_1' b_1', \quad b_2 = p_2' b_2', \quad \ldots \quad b_\mu = p_\mu' b_\mu',}
worin $p_1, p_2 \ldots p_\nu, p_1', p_2' \ldots p_\mu'$ die höchsten Potenzen von $p$ sind, die in $a_1, a_2 \ldots a_\nu, b_1, b_2 \ldots b_\mu$ aufgehen, so kommen alle Primzahlpotenzen $p_1, p_2 \ldots p_\nu$, die grösser als $1$ sind, auch unter den $p_1', p_2' \ldots p_\mu'$ vor, und umgekehrt.}

Um diesen Satz zu beweisen, nehmen wir die Elemente der beiden Basen $A$ und $B$ so geordnet an, dass
\weq{3}{\begin{aligned}
p_1 &\geqq p_2 \geqq \ldots \geqq p_\nu \\
p_1' &\geqq p_2' \geqq \ldots \geqq p_\mu'.
\end{aligned}}
Die in (1) und (2) vorkommenden ganzen Zahlen $a_1', a_2' \ldots a_\nu', b_1', b_2' \ldots b_\mu'$ sind nach ihrer Definition durch $p$ nicht theilbar, und wenn wir also mit $m$ das kleinste gemeinschaftliche Vielfache von $a_1', a_2' \ldots a_\nu'$ bezeichnen, so ist auch $m$ nicht durch $p$ theilbar. Ist aber
\weq{4}{\Theta = A_1^{\alpha_1} A_2^{\alpha_2} \ldots A_\nu^{\alpha_\nu}}
ein beliebiges Element von $S$, so folgt, dass
\[
\Theta^{p_1 m} = 1
\]
sein muss.
